\documentclass[11pt,a4paper]{article}
\usepackage[T1]{fontenc}
\usepackage[utf8]{inputenc}
\usepackage{amsmath,amssymb,amsthm}
\usepackage[margin=1in]{geometry}
\usepackage[colorlinks=true,linkcolor=blue,urlcolor=blue]{hyperref}
\theoremstyle{plain}
\newtheorem{theorem}{Theorem}
\newtheorem{lemma}[theorem]{Lemma}
\newtheorem{proposition}[theorem]{Proposition}
\newtheorem{corollary}[theorem]{Corollary}
\theoremstyle{definition}
\newtheorem{remark}[theorem]{Remark}
% A quoted conjecture can be a single hundred-term formula whose terms are thirty-digit
% numbers. TeX cannot hyphenate those, so without a little stretch every such quote runs off
% the page; the linked-a-file papers are all of that shape.
\setlength{\emergencystretch}{3em}

\title{The empirical closed form for OEIS A223867, proved}
\author{Adrian Perez Fontelles\\ \small Independent researcher}
\date{1 September 2026}
\begin{document}
\maketitle

\begin{abstract}
OEIS A223867 carries an empirical closed form for $a(n)$ --- not a recurrence relating
consecutive terms but an explicit formula. It is true, and it is decidable rather than
empirical. The entry counts arrays subject to a condition on a bounded window of consecutive lines, so the lines of an array are the
vertices of a finite digraph, an array is a walk in it, and $a(n)$ is a walk count on
$S=8960$ vertices. Any function of the form $\sum_j p_j(n)\lambda_j^{\,n}$ satisfies the
monic linear recurrence whose characteristic polynomial is
$\prod_j(x-\lambda_j)^{\deg p_j+1}$; here that polynomial is $q(x)=(x-1)^{25}$, of degree
$25$. Two things are then checked exactly: that the walk count satisfies the same
recurrence from some index on, which the Cayley--Hamilton theorem reduces to a finite
computation, and that the two sequences agree at $25$ consecutive indices past that
point. A monic recurrence determines every later term from its predecessors, so agreement
there is agreement for good.
\end{abstract}

\noindent\small 2020 Mathematics Subject Classification. 05A15, 11B37, 15A18.\normalsize

\section{The conjecture}

OEIS A223867 is ``Number of 4Xn 0..3 arrays with rows and antidiagonals unimodal and columns nondecreasing.''. It has offset $1$ and begins
\[
35, 1225, 24199, 315124, 3017129, 22852913, 144081276, 784071455,\ \dots
\]
The entry states, as a conjecture contributed by its author and never marked settled:
\begin{quote}\raggedright\small
Empirical: a(n) = (1/2906843957821440000)*n\^{}24 + (1/34605285212160000)*n\^{}23 + (2197/1372406261022720000)*n\^{}22 + (37133/608225502044160000)*n\^{}21 + (45197/25276903981056000)*n\^{}20 + (14623277/347557429739520000)*n\^{}19 + (26489599/32011868528640000)*n\^{}18 + (37041601/2667655710720000)*n\^{}17 + (14981899/74392141824000)*n\^{}16 + (19115865583/7532204359680000)*n\^{}15 + (31834111187/1158800670720000)*n\^{}14 + (66986419663/269007298560000)*n\^{}13 + (122476645127999/66283398365184000)*n\^{}12 + (27904576042121/2510734786560000)*n\^{}11 + (949107163427/17557585920000)*n\^{}10 + (100270802691019/470762772480000)*n\^{}9 + (174267558093661/256094948229120)*n\^{}8 + (2011492987640933/1143281018880000)*n\^{}7 + (61760354830949111/16895152834560000)*n\^{}6 + (4216963085471057/703964701440000)*n\^{}5 + (5310598072571189/703964701440000)*n\^{}4 + (2777655817513/391091500800)*n\^{}3 + (19272574007557/3805621142400)*n\^{}2 + (2059892117/1070845776)*n + 1
\end{quote}
As of the ``Last modified'' line on the live entry (Oct 03 2025 22:47:40, revision 8) this is still
recorded as empirical, and nothing on the entry records it as proved. Write
\begin{multline*}
f(n)\;=\;\frac{n^{24}}{2906843957821440000} + \frac{n^{23}}{34605285212160000} \\
+ \frac{2197 n^{22}}{1372406261022720000} + \frac{37133 n^{21}}{608225502044160000} \\
+ \frac{45197 n^{20}}{25276903981056000} + \frac{14623277 n^{19}}{347557429739520000} \\
+ \frac{26489599 n^{18}}{32011868528640000} + \frac{37041601 n^{17}}{2667655710720000} \\
+ \frac{14981899 n^{16}}{74392141824000} + \frac{19115865583 n^{15}}{7532204359680000} \\
+ \frac{31834111187 n^{14}}{1158800670720000} + \frac{66986419663 n^{13}}{269007298560000} \\
+ \frac{122476645127999 n^{12}}{66283398365184000} + \frac{27904576042121 n^{11}}{2510734786560000} \\
+ \frac{949107163427 n^{10}}{17557585920000} + \frac{100270802691019 n^{9}}{470762772480000} \\
+ \frac{174267558093661 n^{8}}{256094948229120} + \frac{2011492987640933 n^{7}}{1143281018880000} \\
+ \frac{61760354830949111 n^{6}}{16895152834560000} + \frac{4216963085471057 n^{5}}{703964701440000} \\
+ \frac{5310598072571189 n^{4}}{703964701440000} + \frac{2777655817513 n^{3}}{391091500800} \\
+ \frac{19272574007557 n^{2}}{3805621142400} + \frac{2059892117 n}{1070845776} \\
+ 1.
\end{multline*}

\section{The count is a walk count}

The entry's defining condition is local: it is a condition on a bounded window of consecutive lines. Take as vertices the
admissible configurations of such a window and put an edge from $u$ to $v$ when $v$ may
follow $u$, which the condition decides by inspection. An admissible array is then exactly a
walk, so with $M$ the adjacency matrix of that digraph on $S=8960$ vertices, and $u,w$ the
vectors recording which windows may start and end an array,
\[
a(n)\;=\;u^{\!\top}M^{\,g(n)}w
\]
for the linear function $g$ that the entry's shape supplies. In particular $a$ is
$C$-finite: it satisfies some monic integer linear recurrence, though not necessarily a short
one.

\section{A closed form is a recurrence}

\begin{lemma}\label{lem:ann}
Let $f(m)=\sum_{j}p_j(m)\lambda_j^{\,m}$ with the $\lambda_j$ distinct and the $p_j$
polynomials, and put $q(x)=\prod_j(x-\lambda_j)^{\deg p_j+1}$. Then $q$ applied to the shift
operator annihilates $f$: writing $q(x)=x^{r}-\sum_{i=1}^{r}c_ix^{r-i}$,
\[
f(m)\;=\;\sum_{i=1}^{r}c_i\,f(m-i)\qquad\text{for every }m.
\]
\end{lemma}

\begin{proof}
The shift operator $E$ acts on $m\mapsto\lambda^{m}p(m)$ as
$(E-\lambda)\bigl(\lambda^{m}p(m)\bigr)=\lambda^{m+1}\bigl(p(m+1)-p(m)\bigr)$, and
$p(m+1)-p(m)$ has degree one less than $p$. So $(E-\lambda)^{\deg p+1}$ kills
$\lambda^{m}p(m)$, and the factors of $q$ commute.
\end{proof}

Here $q(x)=(x-1)^{25}$, of degree $25$; the closed form is a polynomial in $n$, so this is $(x-1)^{25}$, and the recurrence it encodes is
\begin{equation}\label{eq:rec}
a(n)\;=\;25\,a(n-1) - 300\,a(n-2) + 2300\,a(n-3) - 12650\,a(n-4) + 53130\,a(n-5) - 177100\,a(n-6) + 480700\,a(n-7) - 1081575\,a(n-8) + 2042975\,a(n-9) - 3268760\,a(n-10) + 4457400\,a(n-11) - 5200300\,a(n-12) + 5200300\,a(n-13) - 4457400\,a(n-14) + 3268760\,a(n-15) - 2042975\,a(n-16) + 1081575\,a(n-17) - 480700\,a(n-18) + 177100\,a(n-19) - 53130\,a(n-20) + 12650\,a(n-21) - 2300\,a(n-22) + 300\,a(n-23) - 25\,a(n-24) + a(n-25).
\end{equation}

\section{The proof}

\begin{lemma}\label{lem:crit}
With $q$ as above, $a$ satisfies \eqref{eq:rec} for all large $n$ if and only if
$u^{\!\top}M^{\,j}q(M)w=0$ for every $j\ge0$, and it is enough to check $j<S$.
\end{lemma}

\begin{proof}
Substituting the walk expression, the residual $a(n)-\sum_ic_ia(n-i)$ equals
$u^{\!\top}M^{\,g(n)-r}q(M)w$ up to the fixed linear reparametrisation $g$. For the bound,
$M^{S}$ is an integer combination of $I,M,\dots,M^{S-1}$ by Cayley--Hamilton, so the values
$u^{\!\top}M^{j}q(M)w$ satisfy the monic recurrence given by the characteristic polynomial of
$M$; $S$ consecutive zeros therefore force all later ones, and the last nonzero value pins the
threshold exactly.
\end{proof}

\begin{theorem}
$a(n)=f(n)$ for every $n\ge1$.
\end{theorem}

\begin{proof}
The vector $w'=q(M)w$ was formed by $25$ matrix--vector products in exact integer
arithmetic, and $u^{\!\top}M^{j}w'$ evaluated for $j=0,1,\dots$, again exactly. Those integers
vanish from the point corresponding to $n=26$ onwards, and the run was continued until the
working vector was identically zero, which settles every larger $j$ at once. So $a$ satisfies
\eqref{eq:rec} for every $n>25$. By Lemma~\ref{lem:ann} $f$ satisfies \eqref{eq:rec}
identically. Both were then evaluated in exact arithmetic at the $25$ consecutive indices
$n=26,\dots,50$ and agree at each. A monic recurrence of order $25$
determines $a(m)$ from $a(m-1),\dots,a(m-25)$, and for every $m\ge51$ both
$m>25$ and $m-25\ge26$ hold, so induction on $m$ gives $a(m)=f(m)$ for every
$m\ge26$.
Below $n=26$ the recurrence argument gives nothing, since \eqref{eq:rec} is only
established for $a$ from $n=26$ on. The remaining indices $n=1,\dots,25$ were
therefore checked one at a time, directly against the walk count in exact integer arithmetic,
and agree.
\end{proof}

The range is tight in the sense that was checked: the closed form holds from the entry's first term onwards.

\section{Verification}

Four checks, each independent of the algebra above.

First, the walk model reproduces every one of the $18$ terms the entry publishes,
exactly, in integer arithmetic. This is what ties the model to the entry: the model is built
by reading the entry's English, and a misreading gives a different digraph and different
counts.

Second, the closed form was evaluated in exact rational arithmetic --- not floating point ---
at every published term from $n=1$ on, and agrees with the entry's own DATA at each.

Third, the annihilation test of Lemma~\ref{lem:crit} was carried out exactly, and the model
and the closed form were then compared directly at sixty consecutive indices beyond
$n=1$, far past where the proof already settles the matter.

Fourth, the closed form was deliberately perturbed --- by a constant and by a term linear in
$n$ --- and each perturbed form was rejected by the same comparison, so the test is not one
that anything would pass.

\begin{thebibliography}{9}
\bibitem{oeis} The OEIS Foundation, \emph{The On-Line Encyclopedia of Integer
Sequences}, \url{https://oeis.org}, 2026. Sequence A223867.
\bibitem{stanley} R.~P.~Stanley, \emph{Enumerative Combinatorics}, Volume 1, 2nd ed.,
Cambridge University Press, 2011. (Transfer-matrix method, Section 4.7; $C$-finite sequences,
Section 4.1.)
\bibitem{kp} M.~Kauers and P.~Paule, \emph{The Concrete Tetrahedron}, Springer, 2011.
(Closed forms and linear recurrences with constant coefficients.)
\bibitem{horn} R.~A.~Horn and C.~R.~Johnson, \emph{Matrix Analysis}, 2nd ed., Cambridge
University Press, 2013.
\end{thebibliography}

\end{document}
